
\documentclass{article}
\usepackage{colortbl}
\usepackage{makecell}
\usepackage{multirow}
\usepackage{supertabular}

\begin{document}

\newcounter{utterance}

\centering \large Interaction Transcript for game `function\_detective', experiment `list\_active\_pair\_checks', episode 116 with gpt{-}4.1{-}nano.
\vspace{24pt}

{ \footnotesize  \setcounter{utterance}{1}
\setlength{\tabcolsep}{0pt}
\begin{supertabular}{c@{$\;$}|p{.15\linewidth}@{}p{.15\linewidth}p{.15\linewidth}p{.15\linewidth}p{.15\linewidth}p{.15\linewidth}}
    \# & \multicolumn{2}{c}{Player} && \multicolumn{2}{c}{Game Master} \\
    \hline

    \theutterance \stepcounter{utterance}  
    & & & \multicolumn{4}{p{0.6\linewidth}}{
        \cellcolor[rgb]{0.9,0.9,0.9}{
            \makecell[{{p{\linewidth}}}]{
                \texttt{\tiny{[P1$\langle$GM]}}
                \texttt{I have a hidden Python function for you to figure out.} \\
\\ 
\texttt{Your job is to infer what it does and then submit a correct Python implementation.} \\
\\ 
\texttt{Function signature:} \\
\texttt{(items: List{[}int{]}) {-}> List{[}int{]}} \\
\\ 
\texttt{You have 15 turns total.} \\
\texttt{Setup: Active Pair Checks} \\
\\ 
\texttt{You propose function inputs together with a candidate output, and I tell you whether that pair is correct.} \\
\\ 
\texttt{You must make your best SOLVE attempt no later than the second{-}to{-}last turn.} \\
\texttt{Do not wait until the very end.} \\
\\ 
\texttt{You can use these actions:} \\
\texttt{{-} TEST: <inputs>, <candidate\_output>} \\
\texttt{{-} SOLVE: ```python ... ```} \\
\\ 
\texttt{Final answer format:} \\
\texttt{SOLVE: ```python} \\
\texttt{def solution(<value1>):} \\
\texttt{    ...} \\
\texttt{```} \\
\\ 
\texttt{Do not add import statements in your guess.} \\
\texttt{Anything you might need is already available, so just write the function body and return the answer.} \\
\\ 
\texttt{To query graph membership, write exactly one line: `TEST: <input1>, ..., <candidate\_output>`.} \\
\texttt{The last value is the candidate function output; the reply is True iff it matches the hidden function.} \\
\\ 
\texttt{I will reply with OUTPUT: lines only.} \\
\texttt{You must never write OUTPUT: yourself.} \\
\\ 
\texttt{Preloaded examples:} \\
\texttt{1. (({[}5, {-}5, 5{]}, {[}5, {-}5{]}), True)} \\
\\ 
\texttt{Scoring:} \\
\texttt{{-} correctness is required for any non{-}zero score} \\
\texttt{{-} in interactive modes, solving with fewer information requests is better} \\
\\ 
\texttt{Input formatting reminders:} \\
\texttt{{-} strings must be quoted} \\
\texttt{{-} booleans must be True/False} \\
\texttt{{-} use commas between multiple values} \\
\texttt{{-} do not write any extra action tags in the same message} \\
\\ 
\texttt{Format example:} \\
\texttt{You write:} \\
\texttt{TEST: <input1\_number>, <candidate\_output>} \\
\texttt{I answer:} \\
\texttt{OUTPUT: True} \\
\\ 
\texttt{Output only one valid action per turn.} \\
            }
        }
    }
     \\ \\

    \theutterance \stepcounter{utterance}  
    & \multicolumn{4}{p{0.6\linewidth}}{
        \cellcolor[rgb]{0.561,0.737,0.561}{
            \makecell[{{p{\linewidth}}}]{
                \texttt{\tiny{[P1$\rangle$GM]}}
                \texttt{Let's analyze the provided example:} \\
\\ 
\texttt{Preloaded example:} \\
\texttt{{-} Input: ({[}5, {-}5, 5{]}, {[}5, {-}5{]})} \\
\texttt{{-} Output: True} \\
\\ 
\texttt{This suggests that the function might be checking whether the second list is a subset or a specific subset of the first list, or perhaps whether the second list is contained within the first list in some way.} \\
\\ 
\texttt{Given the example:} \\
\texttt{{-} First list: {[}5, {-}5, 5{]}} \\
\texttt{{-} Second list: {[}5, {-}5{]}} \\
\texttt{{-} Result: True} \\
\\ 
\texttt{It seems to be checking if the second list is a subset of the first list.} \\
\\ 
\texttt{Let's test this hypothesis with a few more cases.} \\
\\ 
\texttt{Next, I will test whether the second list is a subset of the first list.} \\
\\ 
\texttt{Turn 1:} \\
\texttt{TEST: {[}1, 2, 3{]}, {[}2, 3{]}, True} \\
\\ 
\texttt{If the function checks subset, it should return True here.} \\
\\ 
\texttt{Please confirm.} \\
            }
        }
    }
    & & \\ \\

    \theutterance \stepcounter{utterance}  
    & & & \multicolumn{4}{p{0.6\linewidth}}{
        \cellcolor[rgb]{0.9,0.9,0.9}{
            \makecell[{{p{\linewidth}}}]{
                \texttt{\tiny{[P1$\langle$GM]}}
                \texttt{Game Violation: ParseError: Invalid format for the current mode.} \\
\texttt{Output ONLY:} \\
\texttt{TEST: ...} \\
\texttt{OR} \\
\texttt{SOLVE: ```python ... ```} \\
            }
        }
    }
     \\ \\

    \theutterance \stepcounter{utterance}  
    & \multicolumn{4}{p{0.6\linewidth}}{
        \cellcolor[rgb]{0.561,0.737,0.561}{
            \makecell[{{p{\linewidth}}}]{
                \texttt{\tiny{[P1$\rangle$GM]}}
                \texttt{TEST: {[}1, 2, 3{]}, {[}2, 3{]}, True} \\
            }
        }
    }
    & & \\ \\

    \theutterance \stepcounter{utterance}  
    & & & \multicolumn{4}{p{0.6\linewidth}}{
        \cellcolor[rgb]{0.9,0.9,0.9}{
            \makecell[{{p{\linewidth}}}]{
                \texttt{\tiny{[P1$\langle$GM]}}
                \texttt{Game Violation: ParseError: Membership probe must contain function inputs followed by a candidate output.} \\
\texttt{Output ONLY:} \\
\texttt{TEST: ...} \\
\texttt{OR} \\
\texttt{SOLVE: ```python ... ```} \\
            }
        }
    }
     \\ \\

    \theutterance \stepcounter{utterance}  
    & \multicolumn{4}{p{0.6\linewidth}}{
        \cellcolor[rgb]{0.561,0.737,0.561}{
            \makecell[{{p{\linewidth}}}]{
                \texttt{\tiny{[P1$\rangle$GM]}}
                \texttt{TEST: {[}1, 2, 3{]}, {[}2, 3{]}, True} \\
            }
        }
    }
    & & \\ \\

    \theutterance \stepcounter{utterance}  
    & & & \multicolumn{4}{p{0.6\linewidth}}{
        \cellcolor[rgb]{0.9,0.9,0.9}{
            \makecell[{{p{\linewidth}}}]{
                \texttt{\tiny{[P1$\langle$GM]}}
                \texttt{Game Violation: ParseError: Membership probe must contain function inputs followed by a candidate output.} \\
\texttt{Output ONLY:} \\
\texttt{TEST: ...} \\
\texttt{OR} \\
\texttt{SOLVE: ```python ... ```} \\
            }
        }
    }
     \\ \\

    \theutterance \stepcounter{utterance}  
    & \multicolumn{4}{p{0.6\linewidth}}{
        \cellcolor[rgb]{0.561,0.737,0.561}{
            \makecell[{{p{\linewidth}}}]{
                \texttt{\tiny{[P1$\rangle$GM]}}
                \texttt{TEST: {[}1, 2, 3{]}, {[}2, 3{]}, True} \\
            }
        }
    }
    & & \\ \\

    \theutterance \stepcounter{utterance}  
    & & & \multicolumn{4}{p{0.6\linewidth}}{
        \cellcolor[rgb]{0.9,0.9,0.9}{
            \makecell[{{p{\linewidth}}}]{
                \texttt{\tiny{[P1$\langle$GM]}}
                \texttt{Game Violation: ParseError: Membership probe must contain function inputs followed by a candidate output.} \\
\texttt{Output ONLY:} \\
\texttt{TEST: ...} \\
\texttt{OR} \\
\texttt{SOLVE: ```python ... ```} \\
            }
        }
    }
     \\ \\

    \theutterance \stepcounter{utterance}  
    & \multicolumn{4}{p{0.6\linewidth}}{
        \cellcolor[rgb]{0.561,0.737,0.561}{
            \makecell[{{p{\linewidth}}}]{
                \texttt{\tiny{[P1$\rangle$GM]}}
                \texttt{TEST: {[}1, 2, 3{]}, {[}2, 3{]}, True} \\
            }
        }
    }
    & & \\ \\

    \theutterance \stepcounter{utterance}  
    & & & \multicolumn{4}{p{0.6\linewidth}}{
        \cellcolor[rgb]{0.9,0.9,0.9}{
            \makecell[{{p{\linewidth}}}]{
                \texttt{\tiny{[P1$\langle$GM]}}
                \texttt{Game Violation: ParseError: Membership probe must contain function inputs followed by a candidate output.} \\
\texttt{Output ONLY:} \\
\texttt{TEST: ...} \\
\texttt{OR} \\
\texttt{SOLVE: ```python ... ```} \\
            }
        }
    }
     \\ \\

    \theutterance \stepcounter{utterance}  
    & \multicolumn{4}{p{0.6\linewidth}}{
        \cellcolor[rgb]{0.561,0.737,0.561}{
            \makecell[{{p{\linewidth}}}]{
                \texttt{\tiny{[P1$\rangle$GM]}}
                \texttt{TEST: {[}1, 2, 3{]}, {[}2, 3{]}, True} \\
            }
        }
    }
    & & \\ \\

    \theutterance \stepcounter{utterance}  
    & & & \multicolumn{4}{p{0.6\linewidth}}{
        \cellcolor[rgb]{0.9,0.9,0.9}{
            \makecell[{{p{\linewidth}}}]{
                \texttt{\tiny{[P1$\langle$GM]}}
                \texttt{Game Violation: ParseError: Membership probe must contain function inputs followed by a candidate output.} \\
\texttt{Output ONLY:} \\
\texttt{TEST: ...} \\
\texttt{OR} \\
\texttt{SOLVE: ```python ... ```} \\
            }
        }
    }
     \\ \\

    \theutterance \stepcounter{utterance}  
    & \multicolumn{4}{p{0.6\linewidth}}{
        \cellcolor[rgb]{0.561,0.737,0.561}{
            \makecell[{{p{\linewidth}}}]{
                \texttt{\tiny{[P1$\rangle$GM]}}
                \texttt{SOLVE: ```python} \\
\texttt{def items(items: list{[}int{]}) {-}> list{[}int{]}:} \\
\texttt{    \# Based on the example, the function seems to check if the second list is a subset of the first list.} \\
\texttt{    \# Since the example confirms this, we'll implement subset check.} \\
\texttt{    \# But the function signature only takes one argument, so perhaps the input is a list of two lists.} \\
\texttt{    \# Let's assume the input is a list of two lists: {[}list1, list2{]}} \\
\texttt{    list1, list2 = items} \\
\texttt{    return list2 if all(elem in list1 for elem in list2) else {[}{]}} \\
\texttt{```} \\
            }
        }
    }
    & & \\ \\

    \theutterance \stepcounter{utterance}  
    & & & \multicolumn{2}{p{0.3\linewidth}}{
        \cellcolor[rgb]{0.9,0.9,0.9}{
            \makecell[{{p{\linewidth}}}]{
                \texttt{\tiny{[GM$|$GM]}}
                \texttt{Game Verdict: LOSS} \\
            }
        }
    }
    & & \\ \\

    \theutterance \stepcounter{utterance}  
    & & & \multicolumn{2}{p{0.3\linewidth}}{
        \cellcolor[rgb]{0.9,0.9,0.9}{
            \makecell[{{p{\linewidth}}}]{
                \texttt{\tiny{[GM$|$GM]}}
                \texttt{Incorrect.} \\
\texttt{Failed Cases:} \\
\texttt{Input: {[}{[}5, {-}5, 5{]}{]} | Expected: {[}5, {-}5{]} | Got: \_\_ERROR\_\_} \\
\texttt{Input: {[}{[}0{]}{]} | Expected: {[}0{]} | Got: \_\_ERROR\_\_} \\
\texttt{Input: {[}{[}1, 1{]}{]} | Expected: {[}1{]} | Got: \_\_ERROR\_\_} \\
            }
        }
    }
    & & \\ \\

    \theutterance \stepcounter{utterance}  
    & & & \multicolumn{2}{p{0.3\linewidth}}{
        \cellcolor[rgb]{0.9,0.9,0.9}{
            \makecell[{{p{\linewidth}}}]{
                \texttt{\tiny{[GM$|$GM]}}
                \texttt{def rmv\_duplicates(items: List{[}int{]}) {-}> List{[}int{]}:} \\
\texttt{    seen = set()} \\
\texttt{    out: List{[}int{]} = {[}{]}} \\
\texttt{    for item in items:} \\
\texttt{        if item not in seen:} \\
\texttt{            seen.add(item)} \\
\texttt{            out.append(item)} \\
\texttt{    return out} \\
            }
        }
    }
    & & \\ \\

    \theutterance \stepcounter{utterance}  
    & & & \multicolumn{2}{p{0.3\linewidth}}{
        \cellcolor[rgb]{0.9,0.9,0.9}{
            \makecell[{{p{\linewidth}}}]{
                \texttt{\tiny{[GM$|$GM]}}
                \texttt{def items(items: list{[}int{]}) {-}> list{[}int{]}:} \\
\texttt{    \# Based on the example, the function seems to check if the second list is a subset of the first list.} \\
\texttt{    \# Since the example confirms this, we'll implement subset check.} \\
\texttt{    \# But the function signature only takes one argument, so perhaps the input is a list of two lists.} \\
\texttt{    \# Let's assume the input is a list of two lists: {[}list1, list2{]}} \\
\texttt{    list1, list2 = items} \\
\texttt{    return list2 if all(elem in list1 for elem in list2) else {[}{]}} \\
            }
        }
    }
    & & \\ \\

    \theutterance \stepcounter{utterance}  
    & & & \multicolumn{2}{p{0.3\linewidth}}{
        \cellcolor[rgb]{0.9,0.9,0.9}{
            \makecell[{{p{\linewidth}}}]{
                \texttt{\tiny{[GM$|$GM]}}
                \texttt{Mode: active\_pair\_checks} \\
\texttt{Accuracy: 0.000} \\
\texttt{Information requests: 0/15} \\
\texttt{Revealed passive examples: 0} \\
\texttt{Efficiency: 1.000} \\
\texttt{Slack: 15} \\
\texttt{Observed{-}pair consistency: 0.000} \\
\texttt{Membership balance: 1 true / 0 false} \\
\texttt{Parse errors: 6} \\
\texttt{Runtime errors: 0} \\
            }
        }
    }
    & & \\ \\

\end{supertabular}
}

\end{document}
